%%% Local macros — all \providecommand to avoid redefinition errors
\providecommand{\Sstab}{S_{\mathrm{stab}}}
\providecommand{\CTB}{\bar{\mathcal{C}}_{\mathrm{TB}}}
\providecommand{\PLVcn}{\mathrm{PLV}_{\mathrm{cn}}}
\providecommand{\tauR}{\tau_{\mathrm{restore}}}
\providecommand{\kappaON}{\kappa_{\mathrm{O_2\text{-}n}}}
\providecommand{\kON}{\kappa_{\mathrm{O_2\text{-}n}}}
\providecommand{\kappaAN}{\kappa_{\mathrm{anaerobic}}}
\providecommand{\calT}{\mathcal{T}}
\providecommand{\calC}{\mathcal{C}}
\providecommand{\calI}{\mathcal{I}}
\providecommand{\calP}{\mathcal{P}}
\providecommand{\Tform}{\tau_{\mathrm{form}}}
\providecommand{\OID}{\mathrm{OID}}
\providecommand{\phiC}{\phi_C}
\providecommand{\phiL}{\phi_L}
\providecommand{\phiTh}{\phi_{\mathrm{th}}}
\providecommand{\omC}{\omega_C}
\providecommand{\omL}{\omega_L}

\chapter{Performance Variance Minimisation in Sprint Running}
\label{ch:variance_min}

\papernote{Source paper: \emph{On the Thermodynamic Consequences of Categorical Flux Propagation Mechanics on Perception: Stability Analysis of Internally-Generated Oscillatory Perturbations on Automatic Motor Substrate During Sprint Running} --- \texttt{parable/docs/physiology/performance-variance-minimisation.tex}}

\begin{chapterabstract}
The foregoing chapters established the physical machinery of biological information processing: the closed charge circuit (Ch.~\ref{ch:musculoskeletal}), autocatalytic metabolic redistribution (Ch.~\ref{ch:metabolism_engine}), functional charge decomposition (Ch.~\ref{ch:charge_quantification}), and the cardiac master oscillator coupled to neural phase dynamics (Ch.~\ref{ch:cardio_neural}). This chapter unifies these components into a single operational framework for measuring individual performance variance during sprint running. The 400\,m sprint provides an ideal experimental window in which experienced runners execute locomotion through cardiac-entrained central pattern generators, freeing higher neural oscillatory function for independent observation. Three quantitative metrics emerge from this architecture---a Stability Index $\Sstab \in [0,1]$, a Thought-Body Coherence measure $\CTB$, and a cardiac-neural Phase-Locking Value $\PLVcn$---linked by the empirical relation $\Sstab = 0.2 + 1.0\,\CTB$ ($R^2 > 0.8$). Clinical stratification across five populations validates the framework as a diagnostic instrument for oscillatory coherence deficits. The chapter also derives the 8000-fold enhancement of neural coupling efficiency attributable to atmospheric O$_2$ paramagnetism and the trans-Planckian temporal precision achievable through gear-ratio cascades across 13 hierarchical oscillatory scales.
\end{chapterabstract}

%% ================================================================
\section{Thought-Motor Dissociation}
%% ================================================================

\begin{multicols}{2}

The central experimental assumption of this chapter is that oscillatory neural thought generation and automatic locomotor output are \emph{dissociable} systems. This must be established before any measurement framework built on their separation can be trusted.

\begin{theorem}[Thought-Motor Dissociation]
\label{thm:ch10-dissociation}
The neural system $\calT$ generating oscillatory thought signatures and the motor system $\mathcal{M}$ producing skeletal output are functionally independent and can operate simultaneously without mutual interference.
\end{theorem}

\begin{proof}[Proof via REM Sleep Physiology]
During REM sleep, subjects exhibit full brainstem muscle atonia: spinal motor neurons are actively inhibited by pontomedullary circuits, producing complete suppression of voluntary movement \citep{chase1990}. Yet throughout this atonia, subjects generate elaborate internally-driven oscillatory narratives indistinguishable subjectively from waking experience---narratives involving complex locomotion, speech, and action sequences. Neural correlates of motor preparation activate in motor cortex, yet no peripheral command reaches skeletal muscle \citep{hobson2002cognitive}.

This constitutes a direct proof: the oscillatory signature-generation apparatus $\calT$ operates at full capacity while motor output $\mathcal{M} \equiv 0$. Therefore $\calT$ and $\mathcal{M}$ are dissociable. \qed
\end{proof}

The dissociation theorem permits the following operational definition.

\begin{definition}[Internal Simulation Operator]
\begin{equation}
\calI : \mathrm{State}_{\mathrm{cog}} \to \mathrm{Experience}_{\mathrm{osc}}
\end{equation}
maps cognitive state (memory, affect, goals) to the oscillatory experiential output. During dreaming $\calI$ is unconstrained; during waking automatic behavior $\calI$ is pegged to external sensory reality through the \emph{reality-pegging operator} $\calP$.
\end{definition}

The resulting waking thought system is:
\begin{equation}
\calT_{\mathrm{wake}} = \calP\!\left(\calI(\mathrm{cognition}),\, \mathcal{R}_{\mathrm{sensory}}\right)
\end{equation}

During an automatic motor task, $\calT_{\mathrm{wake}}$ operates orthogonally to $\mathcal{M}_{\mathrm{auto}}$: thoughts concern strategy, sensation, and affect---not motor commands. This orthogonality is the experimental handle.

\end{multicols}

%% ================================================================
\section{Oscillatory Reality: Three Foundational Pillars}
%% ================================================================

\begin{multicols}{2}

The measurement framework rests on the claim that physical reality is inherently oscillatory. Three independent arguments establish this.

\subsection{Mathematical Necessity}

\begin{theorem}[Oscillatory Recurrence --- Mathematical]
\label{thm:ch10-godel}
Any sufficiently rich formal system that refers to its own states must exhibit periodic recurrence of those states.
\end{theorem}

The argument parallels G\"{o}del self-reference: a system capable of self-modelling generates self-referential loops; self-referential loops in a finite state-space force cyclic trajectories. Every physical dynamical system may be encoded as a formal system over its microstate alphabet; therefore periodic recurrence is a mathematical necessity, not merely a contingent feature.

\subsection{Physical Inevitability}

\begin{theorem}[Oscillatory Recurrence --- Physical]
\label{thm:ch10-poincare}
Every bounded Hamiltonian system with finite phase-space volume returns arbitrarily close to any initial state in finite time.
\end{theorem}

Poincar\'{e}'s recurrence theorem applies directly. Biological subsystems are bounded (finite mass, finite energy budget) and conservative on sufficiently short timescales (adiabatic approximation holds for sub-millisecond neural dynamics). Therefore oscillatory recurrence is physically inevitable \citep{poincare1890}.

\subsection{Thermodynamic Mandate}

Entropy maximisation over a bounded phase space populates all accessible oscillatory modes. For a mode of energy $\hbar\omega$, the thermal occupation is:
\begin{equation}
\langle n_\omega \rangle = \frac{1}{\exp(\hbar\omega / k_B T) - 1}
\end{equation}

At physiological temperature ($k_B T \approx 25\,\mathrm{meV}$), all modes with $\hbar\omega \lesssim k_B T$ have $\langle n_\omega \rangle \gg 1$. Neural oscillatory bands (0.5--100\,Hz) satisfy $\hbar\omega \ll k_B T$ by many orders of magnitude; their population is thermodynamically mandated rather than optionally selected. \emph{Oscillation is the default state of a warm bounded system.}

\end{multicols}

%% ================================================================
\section{Categorical Temporal Emergence}
%% ================================================================

\begin{multicols}{2}

Chapter~\ref{ch:charge_quantification} introduced the Bounded Phase Space $\mathcal{C}(n) = 2n^2$. The present section formalises how \emph{time itself} emerges from the categorical completion of oscillatory events, connecting that earlier analysis to the measurement framework below.

\begin{definition}[Categorical Space with Measure]
A categorical space is a quadruple $(\mathcal{C}, \prec, \mu, \tau)$ where $\mathcal{C}$ is the event set, $\prec$ is a causal order, $\mu : \mathcal{C} \times \mathbb{R}^+ \to [0,1]$ is a completion measure, and $\tau$ is an event topology.
\end{definition}

\begin{axiom}[Irreversibility]
If $\mu(\mathcal{C}, t_1) = 1$ then $\mu(\mathcal{C}, t_2) = 1$ for all $t_2 > t_1$.
\end{axiom}

\begin{theorem}[Temporal Emergence]
\label{thm:ch10-temporal}
For any categorical sequence, a unique temporal coordinate emerges as:
\begin{equation}
T(\mathcal{C}) = \inf\{t : \mu(\mathcal{C}, t) = 1\}
\end{equation}
\end{theorem}

The equivalence $S_{\mathrm{osc}} = S_{\mathrm{cat}}$ established in Ch.~\ref{ch:charge_quantification} (Theorem~\ref{thm:ch09-triple}) means that oscillatory entropy and categorical entropy are the same quantity measured in different frames. The temporal coordinate $T(\mathcal{C})$ is therefore the moment at which an oscillatory mode fully populates its categorical slot---this is precisely the neural synchronisation event that Chapter~\ref{ch:cardio_neural} identified as the coherent processing window $\tau_{\mathrm{sync}} \approx 23\,\mu\mathrm{s}$.

\end{multicols}

%% ================================================================
\section{Biological Maxwell Demon Information Catalysis}
%% ================================================================

\begin{multicols}{2}

Chapter~\ref{ch:metabolism_engine} introduced autocatalytic redistribution as the mechanism by which enzymatic networks channel free energy without violating the Second Law. The BMD formalises this as an information-catalytic process operating at the Landauer limit.

\subsection{Catalytic Selectivity}

An uncatalysed enzymatic channel has prior probability $P_0 \sim 10^{-15}$ of selecting the correct substrate from the accessible microstate pool. A BMD-equipped enzyme achieves posterior probability $P_{\mathrm{BMD}} \approx 10^{-3}$, an enhancement of $10^{12}$ per catalytic event.

The energy cost is bounded by Landauer's erasure principle: selecting one state from $N = 10^{15}$ requires:
\begin{equation}
\Delta F_{\mathrm{Landauer}} = k_B T \ln N \approx 2\,\mathrm{kcal/mol}
\end{equation}
against a free energy budget of $\sim\!30\,\mathrm{kcal/mol}$ per ATP hydrolysis---well within thermodynamic headroom. The BMD therefore operates at $<\!7\%$ of its energy budget on information costs, leaving $>\!93\%$ for mechanical and chemical work.

\subsection{Enzymatic Cascade Architecture}

Multi-step cascades amplify selectivity multiplicatively. Hexokinase provides the canonical example:

\begin{enumerate}
\item Glucose binding ($K_M = 0.1\,\mathrm{mM}$): selectivity $10^5$
\item Phosphate transfer (rate $10^3\,\mathrm{s}^{-1}$): selectivity $10^5$
\item Product release ($K_{\mathrm{cat}} = 100\,\mathrm{s}^{-1}$): selectivity $10^4$
\end{enumerate}

Net cascade enhancement: $\eta_{\mathrm{total}} = 10^{14}$; switching time $\tau_{\mathrm{switch}} = 23\,\mu\mathrm{s}$, consistent with the cardiac-neural coherence window derived in Chapter~\ref{ch:cardio_neural}.

\subsection{Network Synchronisation}

When BMDs operate in synchronised networks, the collective selectivity scales super-additively. For a network of $K$ BMDs with pairwise coupling $\varepsilon$:
\begin{equation}
\eta_{\mathrm{network}} = \eta_{\mathrm{single}}^K \cdot \exp\!\left(\varepsilon K(K-1)/2\right)
\end{equation}

At the cardiac-neural coupling strength $\varepsilon \approx \kON \cdot \tau_{\mathrm{sync}}$ derived in Theorem~\ref{thm:ch09-kappa}, networks of $K \gtrsim 10$ BMDs achieve information bandwidths sufficient for thought-rate processing at $\sim\!3$ events per second.

\end{multicols}

%% ================================================================
\section{S-Entropy Navigation and Five-Dimensional State Space}
%% ================================================================

\begin{multicols}{2}

Chapter~\ref{ch:cardio_neural} introduced the S-entropy coordinate system to characterise the neural phase space. We now make this operational for the running measurement.

\begin{definition}[S-Entropy State Vector]
The S-entropy representation of a neural oscillatory signal assigns five coordinates:
\begin{equation}
\mathbf{s}(t) = (s_K,\, s_T,\, s_S,\, s_C,\, s_I) \in \mathbb{R}^5
\end{equation}
where:
\begin{align}
s_K &= \text{participation ratio (mode distribution)} \\
s_T &= \text{autocorrelation decay timescale} \\
s_S &= \text{sample entropy} \\
s_C &= \text{approach-to-equilibrium rate} \\
s_I &= \text{Shannon entropy of state histogram}
\end{align}
\end{definition}

Navigation through S-space has $\mathcal{O}(1)$ complexity: instead of searching a $10^N$-dimensional microstate space, the system minimises an S-distance:
\begin{equation}
d_S(\mathbf{s}_1, \mathbf{s}_2) = \left[\sum_{j=1}^5 w_j\,(s_{j,1} - s_{j,2})^2\right]^{1/2}
\end{equation}
with weights $w_j$ calibrated to the biological noise floor of each channel. The reduction from exponential to $O(5^2)$ search is the computational dividend of the S-entropy representation.

\subsection{Trans-Planckian Precision via Gear-Ratio Cascades}

Temporal precision appears bounded by the Heisenberg limit:
\begin{equation}
\Delta t \geq \frac{\hbar}{2k_BT} \approx 13\,\mathrm{fs}
\end{equation}

However, gear-ratio coupling between hierarchical oscillatory scales transforms coordinates such that the effective uncertainty becomes:
\begin{equation}
\Delta t_{\mathrm{eff}} = \frac{\Delta t_0}{R_{\mathrm{cascade}}}
\end{equation}

The cascade from cardiac (2.5\,Hz) through neural alpha (10\,Hz), gamma (40\,Hz), cellular kHz, and electronic GHz to molecular THz yields:
\begin{equation}
R_{\mathrm{cascade}} = \underbrace{4}_{\alpha/\mathrm{card}} \times \underbrace{4}_{\gamma/\alpha} \times \underbrace{25}_{\mathrm{cell}/\gamma} \times \underbrace{10^6}_{\mathrm{elec/cell}} \times \underbrace{10^4}_{\mathrm{mol/elec}} = 4\times10^{12}
\end{equation}

From GPS timing ($\Delta t_0 \approx 1\,\mathrm{ms}$):
\begin{equation}
\Delta t_{\mathrm{eff}} \approx \frac{10^{-3}}{4\times10^{12}} \approx 250\,\mathrm{as}
\end{equation}

Combining the cascade with the five-dimensional Fourier enhancement (factor $\sim\!32$) and harmonic network graph topology (factor $\sim\!60$) gives an experimentally measured combined precision of $R_{\mathrm{meas}} = 2847$---consistent with cross-path agreement $\Delta\tau/\tau < 1.3\%$ across three independent measurement pathways \citep{sachikonye2024_variance}.

\end{multicols}

%% ================================================================
\section{Atmospheric Oxygen as Coupling Substrate}
%% ================================================================

\begin{multicols}{2}

The O$_2$ coupling constant $\kON = 4.7\times10^{-3}\,\mathrm{s}^{-1}$ derived in Chapter~\ref{ch:cardio_neural} (Theorem~\ref{thm:ch09-kappa}) deserves mechanistic grounding.

\subsection{Paramagnetic Spin Coupling}

Molecular O$_2$ has triplet ground state $^3\Sigma_g^-$ with two unpaired electrons and net spin $S = 1$, making it uniquely paramagnetic among common atmospheric gases. The Zeeman interaction with endogenous neural magnetic fields ($B_{\mathrm{neural}} \sim 10^{-9}\,\mathrm{T}$) produces spin-state transitions at frequency:
\begin{equation}
\nu_{\mathrm{ESR}} = \frac{g\mu_B B_{\mathrm{neural}}}{h} \approx 28\,\mathrm{Hz}
\end{equation}

This lies squarely in the neural alpha/beta band, providing a spin-resonance channel through which atmospheric O$_2$ couples to cortical oscillatory dynamics.

\subsection{Enhancement Factor}

The ratio of aerobic to anaerobic coupling coefficients is:
\begin{equation}
\frac{\kappaON}{\kappaAN} = \frac{4.7\times10^{-3}}{5.9\times10^{-7}} \approx 8000
\end{equation}

The effective enhancement appearing in BMD efficiency and drug binding efficacy is $\sqrt{8000} \approx 89$.

\subsection{Body-Air Interface During Sprint}

Over a 400\,m sprint ($t \approx 100\,\mathrm{s}$, $v \approx 4\,\mathrm{m/s}$, frontal area $A \approx 0.5\,\mathrm{m}^2$), the runner displaces:
\begin{equation}
V_{\mathrm{air}} = v \cdot A \cdot t = 200\,\mathrm{m}^3 \implies N_{\mathrm{O_2}} \approx 1.1\times10^{27}\,\mathrm{molecules}
\end{equation}

The boundary layer of thickness $\delta \approx 3\,\mathrm{mm}$ maintains $N_{\mathrm{O_2}}^{\mathrm{BL}} \approx 3.4\times10^{22}$ molecules in direct skin contact; 74\% of the neural electromagnetic field penetrates this layer ($B \propto e^{-r/\lambda}$, $\lambda \approx 1\,\mathrm{cm}$). Collectively, the atmospheric O$_2$ reservoir entrains $\sim\!10^{28}$ molecules into coherent coupling with the runner's oscillatory biomechanical dynamics, providing information bandwidth many orders of magnitude above what intracellular O$_2$ alone could supply.

\subsection{Variance Restoration Timescale}

Neural configurational variance $\sigma^2(t)$ obeys exponential relaxation:
\begin{equation}
\frac{d\sigma^2}{dt} = -\frac{\sigma^2}{\tauR}
\end{equation}

with restoration time:
\begin{equation}
\tauR = \frac{1}{\kappaON \cdot N_{\mathrm{O_2}}^{\mathrm{eff}}}
\end{equation}

For $N_{\mathrm{O_2}}^{\mathrm{eff}} \sim 10^{27}$ (atmospheric coupling) this is formally sub-femtosecond; the effective value accounting for neural processing overhead is $\tauR^{\mathrm{eff}} \approx 200\,\mu\mathrm{s}$, matching the $\tau_{\mathrm{switch}}$ of the BMD cascade and the coherent processing window from Ch.~\ref{ch:cardio_neural}. Consequently a complete oscillatory thought signature forms in:
\begin{equation}
t_{\mathrm{thought}} = \tauR^{\mathrm{eff}} \cdot \ln\!\left(\frac{\sigma_0^2}{\sigma_{\mathrm{thresh}}^2}\right) \approx 3.5\,\tauR^{\mathrm{eff}} \approx 200\,\mathrm{ms}
\end{equation}

consistent with empirical estimates of the perceptual latency for a discrete conscious event \citep{libet1983time}.

\end{multicols}

%% ================================================================
\section{The 400\,m Sprint Paradigm}
%% ================================================================

\begin{multicols}{2}

The 400\,m sprint is selected for five independent reasons, each connected to the theoretical architecture developed above.

\paragraph{Motor Automaticity.}
Elite and sub-elite 400\,m runners have automatised their gait through $>10^4$ practice hours. Neuroimaging confirms that during expert locomotion the predominant activation shifts from motor cortex (cortical, voluntary) to cerebellum, brainstem, and spinal cord (subcortical, automatic) \citep{hikosaka2002central}. Central pattern generators (CPGs) in the lumbar spinal cord (L1--L5) sustain the stride cycle without volitional input; the higher neural oscillatory system is therefore genuinely decoupled from locomotor output.

\paragraph{Cardiac Entrainment.}
At 400\,m race pace (heart rate $\sim\!150\,\mathrm{bpm}$, stride rate $\sim\!150\,\mathrm{steps/min}$), cardiac and locomotor cycles entrain at close to 1:1 ratio. Phase-locking value between cardiac and stride phases exceeds $\mathrm{PLV} > 0.7$ in steady-state running \citep{niizeki1993intramuscular}, providing a stable cardiac reference clock against which all other oscillatory signals are measured.

\paragraph{Thought Independence.}
The dissociation theorem (Theorem~\ref{thm:ch10-dissociation}) guarantees that oscillatory thought signatures continue uninterrupted during automatic locomotion. The runner can be instructed to adopt coherent (reality-pegged) or incoherent (reality-divergent) oscillatory patterns; the biomechanical consequences---stability or instability---provide the objective read-out.

\paragraph{Reality Pegging.}
Continuous sensory feedback (proprioception, vestibular, cutaneous) pegs the internal simulation to the actual body state. The Kalman update:
\begin{equation}
\hat{\mathbf{x}}_{k|k} = \hat{\mathbf{x}}_{k|k-1} + \mathbf{K}_k\!\left(\mathbf{y}_k - \mathbf{C}\hat{\mathbf{x}}_{k|k-1}\right)
\end{equation}
quantifies this alignment: when internal prediction $\hat{\mathbf{x}}$ diverges from observed $\mathbf{y}$, the innovation residual is non-zero and stability degrades.

\paragraph{Duration and Statistical Power.}
A 100-second sprint at 200\,ms per thought event yields $\sim\!500$ thought events per run, providing adequate statistical power ($N = 500$) to resolve the three metrics below with confidence intervals of $\pm 0.02$ at $\alpha = 0.05$.

\end{multicols}

%% ================================================================
\section{Central Pattern Generators and Cardiac-Locomotor Coupling}
%% ================================================================

\begin{multicols}{2}

\subsection{CPG Architecture}

The locomotor CPG is a half-centre oscillator with flexor activity $x_F$ and extensor activity $x_E$:
\begin{align}
\frac{dx_F}{dt} &= -x_F + S(w_{FF}x_F - w_{FE}x_E + I_F) \\
\frac{dx_E}{dt} &= -x_E + S(w_{EE}x_E - w_{EF}x_F + I_E)
\end{align}
where $S(z) = (1 + e^{-z})^{-1}$ is the sigmoid, $w_{ij}$ are connection weights with $w_{FE}, w_{EF} > 0$ (mutual inhibition), and $I_i$ are brainstem drive inputs. The cycle period is:
\begin{equation}
T_{\mathrm{cycle}} \approx 2\tau\ln\!\left(\frac{w_{FE}+w_{EF}}{2}\right) \approx 400\,\mathrm{ms}
\end{equation}
consistent with 150 strides/min at race pace.

Proprioceptive feedback from Golgi tendon organs, muscle spindles, and cutaneous foot receptors enters as additive sensory coupling $g_{ij}(\mathbf{s}_{\mathrm{sensory}})$, making the CPG adaptive while preserving its autonomous rhythm.

\subsection{Phase-Locked Cardiac-Locomotor Coupling}

Following the Kuramoto formulation of Chapter~\ref{ch:cardio_neural}, cardiac phase $\phiC$ and locomotor phase $\phiL$ satisfy:
\begin{align}
\frac{d\phiC}{dt} &= \omC + K_{CL}\sin(\phiL - \phiC) \\
\frac{d\phiL}{dt} &= \omL + K_{LC}\sin(\phiC - \phiL)
\end{align}

At 400\,m pace $\omC \approx \omL \approx 2\pi \times 2.5\,\mathrm{Hz}$, which places the system deep inside the 1:1 Arnold tongue. Measured coupling strengths $K_{CL} \approx 0.3\,\mathrm{rad/s}$, $K_{LC} \approx 0.2\,\mathrm{rad/s}$ give a lock-in bandwidth of $\Delta\omega_{\mathrm{lock}} \approx 0.24\,\mathrm{rad/s}$, meaning synchronisation persists across heart rate variations of $\pm 0.04\,\mathrm{Hz}$---robust against mid-race metabolic drift.

\subsection{Thought Phase Against the Cardiac Reference}

With cardiac phase established as the universal reference (see Ch.~\ref{ch:cardio_neural}, Section on cardiac master oscillator), the phase of the $k$-th thought event is:
\begin{equation}
\phiTh^{(k)}(t) = 2\pi \frac{t - t_{R,k}}{T_{RR,k}}
\end{equation}

Phase-locking between thought oscillations and the cardiac reference defines the primary coherence quantity (see Section~\ref{sec:ch10_metrics} below).

\end{multicols}

%% ================================================================
\section{The Three Measurement Metrics}
\label{sec:ch10_metrics}
%% ================================================================

\begin{multicols}{2}

Three scalar quantities characterise the state of the coupled oscillatory system during a 400\,m sprint.

\subsection{Stability Index}

\begin{definition}[Stability Index]
\begin{equation}
\Sstab = 1 - \frac{N_{\mathrm{fall}}^{\mathrm{sim}}}{N_{\mathrm{total}}^{\mathrm{sim}}} \in [0,1]
\end{equation}
where $N_{\mathrm{fall}}^{\mathrm{sim}}$ is the number of stride simulations in which the runner's centre of mass crosses outside the base of support, assessed via a spring-mass biomechanical model calibrated to individual anthropometry.
\end{definition}

$\Sstab = 1$ indicates perfect stride-to-stride stability; $\Sstab = 0$ indicates complete loss of upright posture. The Margin of Stability (MOS):
\begin{equation}
\mathrm{MOS} = x_{\mathrm{BOS}} - x_{\mathrm{COM}} - \frac{v_{\mathrm{COM}}}{\omega_0}
\end{equation}
must remain positive throughout each stance phase; the fraction of strides satisfying this condition defines $\Sstab$ operationally.

\subsection{Thought-Body Coherence}

\begin{definition}[Thought-Body Coherence]
\begin{equation}
\CTB = \frac{1}{N}\sum_{k=1}^{N} \mathcal{C}_k
\end{equation}
where the per-event coherence is:
\begin{equation}
\mathcal{C}_k = \frac{1}{2}\!\left[\cos(\phi_{\mathrm{th},k} - \phiC) + \exp\!\left(-\frac{|f_{\mathrm{th},k} - f_C|}{f_C}\right)\right]
\end{equation}
\end{definition}

The first term measures phase alignment; the second measures frequency alignment. $\CTB \in [0,1]$: $\CTB \to 1$ when thought oscillations are phase- and frequency-locked to the cardiac master; $\CTB \to 0$ when they are fully incoherent.

\subsection{Cardiac-Neural Phase-Locking Value}

\begin{definition}[Cardiac-Neural PLV]
\begin{equation}
\PLVcn = \left|\left\langle e^{i(\phi_{\mathrm{neural}} - \phiC)}\right\rangle_t\right| \in [0,1]
\end{equation}
\end{definition}

This is the circular mean resultant length of the phase-difference distribution, a standard measure of inter-signal synchrony \citep{lachaux1999measuring}. $\PLVcn = 1$ denotes perfect phase-locking; $\PLVcn = 0$ denotes uniformly random phase differences.

\end{multicols}

%% ================================================================
\section{Coherence-Stability Linear Relation}
%% ================================================================

\begin{multicols}{2}

\begin{theorem}[Interface Coherence Theorem]
\label{thm:ch10-coherence}
For automatic motor execution with internal oscillatory simulation operating in parallel, there exists a monotonic relation between $\CTB$ and $\Sstab$:
\begin{equation}
\Sstab = f(\CTB), \quad f' > 0
\end{equation}
\end{theorem}

\begin{proof}[Proof Sketch]
\emph{High coherence} ($\CTB > 0.7$): The internal simulation accurately tracks actual body state. Sensory-predicted state mismatch is small; Kalman innovation residual $\|\mathbf{y}_k - \mathbf{C}\hat{\mathbf{x}}_{k|k-1}\|$ remains below the threshold that would generate corrective motor commands. The automatic CPG substrate runs unperturbed; $\Sstab$ remains high.

\emph{Low coherence} ($\CTB < 0.5$): Internal simulation diverges from actual state. The mismatch generates corrective commands that interfere with the CPG pattern. These perturbations accumulate stride-to-stride, driving the MOS toward zero. $\Sstab$ degrades proportionally.

Linearising the Langevin variance dynamics around steady state and substituting the coherence-dependent potential curvature yields:
\begin{equation}
\boxed{\Sstab = 0.2 + 1.0 \cdot \CTB}
\end{equation}
with intercept 0.2 representing the residual stability of the purely automatic (coherence-independent) CPG substrate, and slope 1.0 reflecting unit coupling between thought-body phase alignment and stride-to-stride biomechanical stability. Empirically $R^2 > 0.8$ \citep{sachikonye2024_variance}. \qed
\end{proof}

The intercept $\beta_0 = 0.2$ is experimentally testable: it is the stability of a runner with completely incoherent oscillatory thought ($\CTB \approx 0$), which should approach but not reach zero because the CPG substrate maintains baseline stride continuity independent of higher oscillatory state.

\end{multicols}

%% ================================================================
\section{Clinical Stratification}
%% ================================================================

\begin{multicols}{2}

The three metrics provide differential diagnostic power across populations. Table~\ref{tab:ch10_clinical} shows reference values from five groups measured under the sprint protocol.

\end{multicols}

\begin{table}[!htbp]
\centering
\caption{Reference values of the three oscillatory coherence metrics across five clinical and performance groups measured during the 400\,m sprint protocol. Values are mean $\pm$ SD. Healthy: general population athletes. MDD: major depressive disorder. All groups $n \geq 20$.}
\label{tab:ch10_clinical}
\begin{tabular}{lcccc}
\hline
\textbf{Group} & $\Sstab$ & $\CTB$ & $\PLVcn$ \\
\hline
Meditation experts & $0.99 \pm 0.01$ & $0.91 \pm 0.04$ & $0.87 \pm 0.05$ \\
Healthy athletes & $0.97 \pm 0.02$ & $0.82 \pm 0.06$ & $0.71 \pm 0.07$ \\
Anxiety disorder & $0.73 \pm 0.08$ & $0.61 \pm 0.09$ & $0.54 \pm 0.10$ \\
Major depression & $0.81 \pm 0.07$ & $0.67 \pm 0.08$ & $0.48 \pm 0.09$ \\
Schizophrenia & $0.52 \pm 0.12$ & $0.43 \pm 0.11$ & $0.36 \pm 0.11$ \\
\hline
\end{tabular}
\end{table}

\begin{multicols}{2}

Several features of Table~\ref{tab:ch10_clinical} merit comment.

\paragraph{Meditation vs.\ healthy.}
Meditation experts show $\Sstab \approx 0.99$, approaching the theoretical ceiling. The mechanism is enhanced cardiac-neural entrainment from long-term attentional training, which tightens the feedback loop in the reality-pegging operator $\calP$ and keeps $\CTB$ systematically above 0.9.

\paragraph{Anxiety vs.\ depression.}
Anxiety yields lower $\Sstab$ than depression ($0.73$ vs.\ $0.81$) despite higher $\PLVcn$ ($0.54$ vs.\ $0.48$). This inversion reflects the distinct mechanisms: anxiety generates high-frequency oscillatory interference that disrupts CPG stride timing (low $\Sstab$) even when cardiac-neural phase-locking is partially maintained; depression produces global oscillatory slowing that reduces PLV without generating the disruptive interference.

\paragraph{Schizophrenia.}
The lowest values across all three metrics ($\Sstab = 0.52$, $\CTB = 0.43$, $\PLVcn = 0.36$) indicate profound disruption of the cardiac-neural coupling architecture, consistent with the coherence deficit picture developed in Chapter~\ref{ch:cardio_neural} (Section on disease as coherence deficit). The stability floor at $\Sstab = 0.52 > 0.2$ confirms the CPG substrate remains partially functional even when higher oscillatory coherence is severely degraded.

\paragraph{Therapeutic floor.}
The intercept $\beta_0 = 0.2$ sets a minimum stability achievable by any individual with a functioning CPG, regardless of oscillatory state. This is the \emph{therapeutic floor}: pharmacological or neuromodulatory interventions that restore $\CTB$ above 0.5 immediately lift $\Sstab$ above $0.7$, the threshold compatible with independent ambulation.

\end{multicols}

%% ================================================================
\section{Psychon Geometry: Oscillatory Thought Quanta}
%% ================================================================

\begin{multicols}{2}

The oscillatory thought event has a definite physical structure. We call the minimal complete thought unit the \emph{oscillatory thought quantum} (OTQ), defined by:

\begin{enumerate}
\item A molecular geometry $\mathbf{G} = \{\mathbf{r}_1,\ldots,\mathbf{r}_N\}$ of $N \sim 10^3$--$10^4$ O$_2$ molecules distributed in concentric shells around a localised electron vacancy (oscillatory hole)
\item An oscillatory signature $\boldsymbol{\omega} = (\omega_1,\ldots,\omega_M)$ phase-locked to the cardiac reference
\item S-entropy coordinates $\mathbf{s} \in \mathbb{R}^5$ uniquely identifying the thought type
\item Formation time $\Tform \approx 200\,\mathrm{ms}$
\item Stability lifetime $\sim\!2$--$3\,\mathrm{s}$ (specious present duration)
\end{enumerate}

\paragraph{Information content.}
The OTQ carries approximately 50 bits of information---the log$_2$ of the equivalence class of geometrically equivalent configurations $|\{[\psi]\}| \approx 10^{15}$. BMD cascade filtering reduces this space to a single actual geometry per thought event at an energy cost of $\sim\!3$ ATP molecules ($\approx 1.5\times10^{-19}\,\mathrm{J}$).

\paragraph{Shell structure.}
O$_2$ molecules organise in four concentric shells around the oscillatory hole:
\begin{itemize}
\item Inner shell (1--5\,\AA, $N_1 \approx 20$): direct electrostatic stabilisation
\item Middle shell (5--10\,\AA, $N_2 \approx 60$): indirect coupling
\item Outer shell (10--20\,\AA, $N_3 \approx 200$): long-range correlation
\item Halo (20--50\,\AA, $N_4 \approx 10^3$): weak background coupling
\end{itemize}

The total $\sim\!1300$ molecules per OTQ constitute the physical substrate of a single oscillatory thought event.

\paragraph{Variance dynamics.}
Configurational variance $\sigma^2(t) = N^{-1}\sum_i\|\mathbf{r}_i(t) - \langle\mathbf{r}_i\rangle\|^2$ obeys:
\begin{equation}
\sigma^2(t) = \sigma_{\mathrm{ss}}^2 + (\sigma_0^2 - \sigma_{\mathrm{ss}}^2)\,e^{-t/\tauR}
\end{equation}

Without BMD activation the steady-state variance diverges (flat potential, no localisation). With BMD the potential becomes sharply convex, reducing $\sigma_{\mathrm{ss}}$ to $\sim\!0.5\,\text{\AA}$---below the oscillatory hole stabilisation threshold---and thereby freezing the geometry into a single OTQ.

\end{multicols}

%% ================================================================
\section{Chapter Summary}
%% ================================================================

\begin{multicols}{2}

This chapter has synthesised the physical substrates established in Chapters~\ref{ch:musculoskeletal}--\ref{ch:cardio_neural} into a unified measurement framework for individual oscillatory performance during sprint running. The main results are as follows.

\paragraph{Dissociation theorem.}
REM sleep physiology proves that oscillatory thought generation ($\calT$) and motor output ($\mathcal{M}$) are functionally dissociable. During automatic locomotion in experienced runners, $\calT$ operates orthogonally to the CPG, permitting independent measurement.

\paragraph{Atmospheric O$_2$ enhancement.}
Paramagnetic coupling coefficient $\kappaON = 4.7\times10^{-3}\,\mathrm{s}^{-1}$ exceeds the anaerobic baseline by a factor of 8000, with an effective BMD efficiency enhancement of $\sqrt{8000}\approx 89$. The 400\,m sprint entrains $\sim\!10^{28}$ atmospheric O$_2$ molecules into coherent coupling, supplying information bandwidth orders of magnitude above intracellular O$_2$ alone.

\paragraph{Trans-Planckian precision.}
Gear-ratio cascades across 13 hierarchical oscillatory scales (GPS $\to$ cardiac $\to$ neural $\to$ molecular $\to$ nuclear) reduce effective temporal uncertainty from the GPS baseline of 1\,ms to attosecond resolution in S-entropy transformed coordinates. Experimentally measured cascade factor $R_{\mathrm{meas}} = 2847$; cross-pathway consistency $\Delta\tau/\tau < 1.3\%$.

\paragraph{Three metrics.}
The Stability Index $\Sstab$, Thought-Body Coherence $\CTB$, and Cardiac-Neural PLV $\PLVcn$ collectively characterise the oscillatory coupling state. They are linked by:
\begin{equation}
\Sstab = 0.2 + 1.0\,\CTB \quad (R^2 > 0.8)
\end{equation}
The intercept 0.2 is the CPG-autonomous stability floor; the slope 1.0 reflects unit coupling between phase alignment and biomechanical stability.

\paragraph{Clinical stratification.}
Reference values stratify cleanly across five populations: meditation experts ($\Sstab = 0.99$, $\CTB = 0.91$, $\PLVcn = 0.87$), healthy athletes ($0.97$, $0.82$, $0.71$), anxiety disorder ($0.73$, $0.61$, $0.54$), major depression ($0.81$, $0.67$, $0.48$), and schizophrenia ($0.52$, $0.43$, $0.36$). The therapeutic floor at $\CTB \approx 0.5$ corresponds to the $\Sstab = 0.7$ threshold compatible with independent ambulation---providing a quantitative target for neuromodulatory intervention.

Part~III of the monograph is hereby complete. The physical substrate of oscillatory neural function has been derived from the closed charge circuit (Ch.~\ref{ch:musculoskeletal}) through autocatalytic metabolism (Ch.~\ref{ch:metabolism_engine}), functional charge decomposition and the Bounded Phase Space $\mathcal{C}(n) = 2n^2$ (Ch.~\ref{ch:charge_quantification}), the cardiac master oscillator and its neural coupling constant $\kON$ (Ch.~\ref{ch:cardio_neural}), and the unified variance-minimisation framework for the running individual (this chapter). Part~IV will extend this physical substrate to the explicit mathematics of oscillatory coherence as the categorical completion of bounded dynamical systems.

\end{multicols}
